\documentclass[10pt,a4paper]{article}
\usepackage[T1]{fontenc}
\usepackage[utf8]{inputenc}
\usepackage{lmodern}
\usepackage[margin=2cm]{geometry}
\usepackage{xcolor}
\usepackage{graphicx}
\usepackage{enumitem}
\usepackage{titlesec}
\usepackage{xurl}
\usepackage[hidelinks]{hyperref}
\usepackage{microtype}

\definecolor{linkblue}{RGB}{20,60,140}
\hypersetup{colorlinks=true, urlcolor=linkblue, linkcolor=linkblue}
\urlstyle{same}
\pagestyle{plain}
\setlength{\parindent}{0pt}

\titleformat{\section}{\large\scshape}{}{0em}{}[\vspace{-0.6ex}\titlerule]
\titlespacing*{\section}{0pt}{2.2ex}{1.2ex}
\titleformat{\subsection}{\normalsize\itshape}{}{0em}{}
\titlespacing*{\subsection}{0pt}{1.4ex}{0.6ex}

\newlist{publist}{enumerate}{1}
\setlist[publist]{leftmargin=2.6em, labelsep=0.6em, itemsep=0.5ex,
                  topsep=0.3ex, align=left}

% \entry{title}{date}: title line with an optional right-aligned date
\newcommand{\entry}[2]{\par\addvspace{0.9ex}%
  \begin{minipage}[t]{\dimexpr\linewidth-3cm}\raggedright #1\strut\end{minipage}%
  \hfill\begin{minipage}[t]{2.8cm}\raggedleft #2\end{minipage}\par}
\newlist{sublist}{itemize}{1}
\setlist[sublist]{label=\textbullet, leftmargin=2em, labelsep=0.6em, itemsep=0.2ex, parsep=0pt,
                  topsep=0.4ex, font=\small, before=\small}

% summary: indented, justified paragraphs under an entry
\newenvironment{summary}{\par\addvspace{0.6ex}\small
  \begin{list}{}{\leftmargin=1em\rightmargin=0pt\topsep=0pt
  \parsep=0.5ex\itemsep=0pt}\item\relax}{\end{list}}
% \meta{text}{date}: secondary line under an entry
\newcommand{\meta}[2]{\nopagebreak{\small\color{black!75}%
  \if\relax\detokenize{#2}\relax #1\else
  \begin{minipage}[t]{\dimexpr\linewidth-3cm}\raggedright #1\strut\end{minipage}%
  \hfill\begin{minipage}[t]{2.8cm}\raggedleft #2\end{minipage}\fi\par}}

\begin{document}

\begin{center}
{\LARGE\scshape George Marchment}\\[1.2ex]
PhD candidate, Paris-Saclay University, LISN Laboratory\\[0.4ex]
\href{mailto:george.marchment@universite-paris-saclay.fr}{george.marchment[at]universite-paris-saclay.fr}\\[0.4ex]
\href{https://scholar.google.com/citations?user=KqPx4dkAAAAJ&hl=fr&oi=ao}{Google Scholar} \enspace$\cdot$\enspace \href{https://github.com/George-Marchment}{GitHub} \enspace$\cdot$\enspace \href{https://linkedin.com/in/george-marchment-9900ba1a9}{Linkedin} \enspace$\cdot$\enspace \href{https://orcid.org/0000-0002-4565-3940}{ORCID}
\end{center}

\medskip
I am a 3rd year PhD candidate studying computer science at Paris-Saclay University. My research interests include data science, graph algorithms, visualisation, AI research and bioinformatics.

\section*{Education}
\entry{\textbf{PhD Candidate}}{2023--present}
\meta{Paris-Saclay University}{}
\meta{Thesis: \emph{A graph-based framework for scientific workflows: modelling, abstraction, and visualisation for reuse}}{}
\meta{Supervisors: \href{https://sarah.cohen-boulakia.eu/}{Sarah Cohen-Boulakia} (Paris-Saclay University) and \href{https://research.pasteur.fr/fr/member/frederic-lemoine/}{Frédéric Lemoine} (Pasteur Institute)}{}
\begin{summary}
\textbf{Modern research has become increasingly data-intensive}, with fields such as bioinformatics generating vast amounts of raw data. Turning this data into meaningful results relies on large, complex analysis pipelines, making effective data management and integration central to scientific discovery. At the same time, growing awareness of the reproducibility crisis has highlighted the need to build these pipelines in a more reproducible way, since science advances by reproducing and reusing prior work.

Scientific workflows, and in particular code-based systems such as Nextflow and Snakemake, have emerged as a solution to provide better reproducibility, offering software encapsulation, scalability and flexibility. Reusing workflows lets researchers build on existing components and supports \textbf{collaborative and interdisciplinary research}. However, workflows are often complex and defined in heterogeneous ways, which makes them hard to share, understand and discover, ultimately limiting their reuse.

To provide an element of solution to this problem, this PhD work proposes a framework that: (i) models heterogeneous workflows as a \textbf{homogeneous graph structure}, bridging the implicit structures of code-based implementations and explicit workflow representations; (ii) formalises graph-based methods to automatically \textbf{generate simplified representations of workflows at multiple levels of granularity}, reducing their perceived complexity; and (iii) applies \textbf{information visualisation} principles to present these structures with a \textbf{consistent visual metaphor}, enabling interactive exploration across simplification levels and linking visuals directly to the source code. These solutions have been implemented in multiple open-source tools, making them directly available to the research community.

Together, these contributions, at the intersection of \textbf{data management and integration}, \textbf{visualisation} and \textbf{open-source software development}, aim to support the overarching goal of enhancing workflow transparency and understandability, in order to facilitate workflow reuse.
\end{summary}

\entry{\textbf{Master's Degree in Bioinformatics, Computational Biology Track, with First Class Honours}}{2021--2023}
\meta{\href{https://digicosme.cnrs.fr/labex-digicosme-bourses-master/}{Digicosme scholarship for excellence}}{}
\meta{Paris-Saclay University}{}
\begin{sublist}
  \item \textbf{Advanced Databases:} query optimisation and rewriting, indexing structures, query plans, schema optimisation and data storage strategies.
  \item \textbf{Massive Data Project:} designed and built a web application to query and enrich a large database. Covered client/server architecture, database design for data volume, security awareness and user needs.
  \item \textbf{Omics Data Analysis and Integrative Biology:} statistical modelling and integration of heterogeneous high-throughput data.
  \item \textbf{Models for Machine Learning:} supervised regression and classification, cross-validation, model selection, random forests, LASSO, SVM.
  \item \textbf{Advanced Machine Learning:} semi-supervised learning, ensemble learning, out-of-distribution learning, transfer learning. Included a group project reproducing a recently published paper and writing a critical report.
\end{sublist}

\entry{\textbf{Double Bachelor's Degree in Mathematics and Computer Science}}{2018--2021}
\meta{Université Paris-Saclay}{}

\section*{Publications}
\subsection*{Peer-reviewed international journals}
\begin{publist}[label={[J\arabic*]}]
  \item \textbf{George Marchment}, Sarah Cohen-Boulakia, Alain Denise and Frédéric Lemoine. Workflow abstractions for automatic modularity: improving workflow code reusability. \emph{Future Generation Computer Systems}, 187:108776, February 2027. \url{https://doi.org/10.1016/j.future.2026.108776} (PDF: \url{https://hal.science/hal-05737452})
  \item \textbf{George Marchment}, Bryan Brancotte, Marie Schmit, Frédéric Lemoine, and Sarah Cohen-Boulakia. BioFlow-Insight: facilitating reuse of Nextflow workflows with structure reconstruction and visualization. \emph{NAR Genomics and Bioinformatics}, 6(3):lqae092, September 2024. \url{https://doi.org/10.1093/nargab/lqae092}
  \item Marine Djaffardjy, \textbf{George Marchment}, Clémence Sebe, Raphael Blanchet, Khalid Belhajjame, Alban Gaignard, Frédéric Lemoine, and Sarah Cohen-Boulakia. Developing and reusing bioinformatics data analysis pipelines using scientific workflow systems. \emph{Computational and Structural Biotechnology Journal}, 21:2075--2085, January 2023. \url{https://doi.org/10.1016/j.csbj.2023.03.003}
\end{publist}

\subsection*{Peer-reviewed international conferences}
\begin{publist}[label={[C\arabic*]}]
  \item Lénora Buggenhoudt, \textbf{George Marchment}, Anne-Flore Cabouat, Frédéric Lemoine, Caroline Appert, and Sarah Cohen-Boulakia. MetroFlow: Visualizing Bioinformatics Workflow Source Code as Interactive Metro-Maps. \emph{IEEE Working Conference on Software Visualization (VISSOFT)}, September 2026. \url{https://hal.science/hal-05737427v1}
  \item Marie Schmit, Ulysse Le Clanche, \textbf{George Marchment}, Sarah Cohen-Boulakia, Olivier Dameron, Alban Gaignard, Frédéric Lemoine, and Hervé Ménager. A Standards-Based Knowledge Graph that Bridges Scientific Workflows, Run-Time Provenance, and Tool Registries. \emph{Semantic Web Applications and Tools for Health Care and Life Sciences}, March 2026. \url{https://hal.science/hal-05517640/}
\end{publist}

\subsection*{Posters}
\begin{publist}[label={[P\arabic*]}]
  \item \textbf{George Marchment}, Bryan Brancotte, Marie Schmit, Frédéric Lemoine, and Sarah Cohen-Boulakia. (2024) BioFlow-Insight: Leveraging Graph representations to Facilitate the Reuse of Bioinformatics Workflows. Poster presented at the \emph{ACM REP} Conference, Rennes, France. \url{https://gitlab.liris.cnrs.fr/sharefair/posters/acm-rep-2024/-/blob/main/Poster\%20ACM\%20REP.pdf?ref_type=heads}
  \item \textbf{George Marchment}, Marie Schmit, Clémence Sebe, Frédéric Lemoine, Hervé Ménager, and Sarah Cohen-Boulakia. (2024) Representing bioinformatics Nextflow workflows in RO-Crate: challenges and opportunities. Poster presented at the \emph{Semantic Web Applications and Tools for Health Care and Life Science} Conference, Leiden, The Netherlands. \url{https://hal.science/hal-04540040}
\end{publist}

\subsection*{Tutorials given}
\entry{\textbf{Computational Reproducibility With Scientific Workflows: Analysing viral genomes with Nextflow}}{July 2025}
\meta{ACM REP \enspace$\cdot$\enspace Tutorial: \url{https://george-marchment.github.io/acmrep25/}}{}

\subsection*{Book chapters}
\begin{publist}[label={[B\arabic*]}]
  \item Sarah Cohen-Boulakia, Frédéric Lemoine, \textbf{George Marchment}, Marine Djaffardjy, Alban Gaignard, Clémence Sebe, and Khalid Belhajjame. Reuse and Reproduce Bioinformatic Pipelines using Scientific Workflow Systems. In Berlin Universities Publishing, editor, \emph{Workflow systems for large-scale scientific data analysis}. Berlin Universities Publishing, 2026. \url{https://pasteur.hal.science/pasteur-05640482/}
  \item Sarah Cohen-Boulakia, \textbf{George Marchment}, Thomas Cokelaer and Frédéric Lemoine. Reprohackathons: Training efforts to increase Bioinformatics Reproducibility using Scientific Workflow systems. In Berlin Universities Publishing, editor, \emph{Workflow systems for large-scale scientific data analysis}. Berlin Universities Publishing, 2026. \url{https://pasteur.hal.science/pasteur-05640478v1}
\end{publist}

\section*{Talks}
\entry{A graph-based framework for scientific workflows: modelling, abstraction, and visualisation for reuse (Planned)}{November 2026}
\meta{\emph{EMBL} Scientific Workflows Club}{}

\entry{Facilitating reuse of Nextflow workflows with structure reconstruction and metro-map visualisation with BioFlow-Insight and MetroFlow}{September 2026}
\meta{FAIR metadata for bioinformatics software workshop at \emph{ECCB 2026}}{}

\entry{Workflow abstractions to improve understandability and reusability of workflow code}{April 2026}
\meta{Visiting the \emph{VIDA Center} (New York University)}{}

\entry{MetroFlow: Generating Interactive Metro-Maps from Workflow Code}{}
\meta{ShareFAIR annual meeting (Pasteur Institute)}{December 2025}
\meta{Nextflow Summit \enspace$\cdot$\enspace Talk: \url{https://summit.nextflow.io/2025/virtual/agenda/metroflow-generating-interactive-metro-maps-from-workflow-code/}}{October 2025}

\entry{Facilitating the reuse of scientific workflows with structure abstraction}{June 2025}
\meta{Visiting the \emph{Software Sustainability Institute} (University of Manchester)}{}

\entry{BioFlow-Insight: facilitating reuse of Nextflow workflows with structure reconstruction and visualisation}{}
\meta{ShareFAIR annual meeting (Pasteur Institute)}{November 2024}
\meta{Team seminar (Paris-Saclay University, LISN)}{January 2025}

\entry{Initial findings on workflow process similarity}{November 2023}
\meta{ShareFAIR kickoff (Paris-Saclay University)}{}

\section*{Software}
\entry{\textbf{MetroFlow}}{}
\meta{Website: \url{https://metroflow.pasteur.cloud/} \enspace$\cdot$\enspace Source code: \url{https://gitlab.pasteur.fr/sharefair/metroflow}}{}

\entry{\textbf{BioFlow-Insight}}{}
\meta{Website: \url{https://bioflow-insight.pasteur.cloud/} \enspace$\cdot$\enspace Source code: \url{https://gitlab.liris.cnrs.fr/sharefair/bioflow-insight}}{}

\entry{\textbf{Github-Crawler}}{}
\meta{Source code: \url{https://gitlab.liris.cnrs.fr/sharefair/github-crawler}}{}

\end{document}
